\scanpage{66}
syst. \inserted{d'isom} en question serait \unsure{invarié|inversé} par ledit \unsure{autom},
absurde. Ceci montre qu'il ne peut y avoir d'unicité quand il existe une ps.droite qui est
un bigône (pas non plus bien sûr s'il y a
un monogone, i.e. s'il s'agit d'un faisceau
de ps.dr) -- mais cela signifie \unsure{just\textsuperscript{mt}|justement} qu'il
s'agit d'un double faisceau \unsure{cont\textsuperscript{nt}} une ps.droite qui joint les deux foyers\ldots

On peut maint\textsuperscript{nt}, dans le cas où il n'y a
ni monogône ni bigône, définir pour tt sommet $s$ $\Omega(s)$ comme l'ens. \unsure{quot\textsuperscript{nt}|quotient} de l'ens
des $\Omega(D)$ ($D$ ps. dr ne passant pas par $s$)
par le syst des isoms $\varphi_{s;D,D'}$. C'est un ens
à deux éléments, \deleted{\unsure{pour}} $s$ variable il faut en
faire un syst. local sur $K^1$ i.e. pour
tt segment $\unsure{I'|I}$ de $K^1$ d'extrémités $s,s'$, il faut
définir une bijection canonique $\Omega(s)\simeq\Omega(s')$

Par hyp. il existe une droite $D$ qui ne passe
par $s$ ni $s'$, et \deleted{par} on définit $\Omega(s)\simeq\Omega(s')$ comme
le composé
\diagram{\begin{tikzcd}[ampersand replacement=\&,column sep=small,row sep=small]
\Omega(s)\arrow[rr,"\sim"]\arrow[dr,"\sim"'] \& \& \Omega(s')\\
\& \Omega(D)\arrow[ur] \&
\end{tikzcd}}
, il résulte de b) que ça ne dépend pas du choix de $D$
% archive sheet number: 65
